\chapter{Physics motivation}
TO BE DECIDED
%PUT THE FOLLOWING BETWEEN THE ELECTROMAGNETIC AND NUCLEAR INTERACTION SECTIONS!
%The electromagnetic interactions described in the preceding section are well understood from both theoretical and experimental perspectives. Nuclear interactions are also a crucial aspect of transport, but are not nearly as well understood from the theoretical perspective. Interactions between a projectile and atoms of the target are ultimately described by quantum electrodynamics (QED), the most accurate physical theory yet devised; by contrast, nuclear interactions are many-body problems that defy present-day calculational methods at the most fundamental level. That is, the particles that participate in nuclear interactions are themselves composites (nuclei contain nucleons, and nucleons contain quarks and gluons), and the fundamental theory that describes these interactions is quantum chromodynamics (QCD), which is only tractable in the limit of interactions with large momentum transfers. QCD has not yet been successfully applied to nucleus–nucleus collisions at the energies of interest here. The lack of a fundamental theory has led to the development of many semi-empirical models to describe nucleus–nucleus interactions, and considerable effort continues to be put into development of these models and benchmarking them (10) against the limited set of pertinent data that are available (2).
%https://www.frontiersin.org/journals/oncology/articles/10.3389/fonc.2016.00065/full

%YOU CAN ADD A DEDICATED SECTION OR SUBSECTION WHICH DESCRIBES THE IMPORTANCE OF CROSS SECTION MEASUREMENT IN PS AND SRP. the information is in Are Further Cross Section Measurements Necessary for Space Radiation Protection or Ion Therapy Applications? Helium Projectiles paper.
% here you have to write (citing the FOOT CDR AT PAGE 8) that studies [reference 17] have shown that scaling by 20% the cross section for inealstic processes in a Monte Carlo transport code (SHIELDIT) modified the RBE calculation by 3%. Here write also the reason why we need such stringent reqiurements on the cross section: such stringnt requirements will translate to stronger constraints in TPS and allow them to take a leap from conservative treatments ("RANGE VERIFICATION PAPER") to predictive ones. here you can cite the 3-5% requirement of dose delivery uncertainty on the paper "https://www.tandfonline.com/doi/full/10.1080/0284186X.2016.1246801%" on my desktop.

% IN "RANGE VERIFICATION" paper there is the citation that conservative treatments are used nowadays because of the uncertainties arising from fragment unwanted energy deposits